\section{One probability space for the different laws}\label{R:sec:renewal}
The elementary renewal identities must precede the changes of measure. They concern complete irreducible pieces. In particular, no independence between the length, height and displacement of one piece is assumed.

\subsection{Ports, strips and the finite input}
Fix one port on a horizontal row boundary. A strict bridge of height $h\in\N$ has its other port on the boundary $h$ rows above and all its vertices in the open strip between these boundaries. One row has Euclidean height $d_0=\sqrt3/2$. Write $\mathcal B_h$ for this set, with the terminal port unrestricted, and set
\[
 B_h=\sum_{\gamma\in\mathcal B_h}\rho^{L(\gamma)},\qquad
 M_h=\sum_{\gamma\in\mathcal B_h}L(\gamma)\rho^{L(\gamma)},\qquad B_0=1.
\]
The term at height zero denotes an empty concatenation, not a nonempty walk. For every port path its diameter includes both ports as well as its visited vertices. Heights and transverse displacements use the actual staggered honeycomb port lattice. All paths and all translations below retain their starting port; there is never an extra sum over its position.

Here is the positive finite input for our main route. Its proof belongs to the uniform marked-polygon calculation and its exterior sewing argument.
\begin{theorem}[Finite bridge estimates]\label{R:renew:finite-input}
There are constants $c,C>0$ such that, for every sufficiently large integer $h$,
\[
 c h^{-1/4}\le B_h\le C h^{-1/4},\qquad M_h\le C h^{13/12},\qquad
 B_h\{L\ge c h^{4/3}\}\ge c h^{-1/4}.
\]
For all $h\ge1$ and $x\ge1$, one also has
\[
 B_h\{D>x\}\le C x^{-1/4}.
\]
The same estimates hold in the symmetry-related pure directions.
\end{theorem}
These are the bounded-factor estimates of \cite[Theorem~1.1]{companionU}. All finitely many smaller-height masses and first moments are finite as well: append a fixed straight bridge to inject their paths into a sufficiently tall strip, with bounded weight and additive length changes. We enlarge constants to include these heights. The first moment and the positive mass of long bridges play different roles: the former bounds a truncated irreducible reward, whereas the latter supplies a lower length deficit. Exponent-only versions lead to the weaker, explicitly stated variants later in the paper.

\subsection{Factorization and marked rewards}
An intermediate row boundary crossed exactly once is a renewal boundary. Cutting a strict bridge at every renewal boundary gives an ordered list of irreducible bridges, each of positive integral height. Conversely, a list of irreducibles concatenates to a strict bridge: its pieces occupy disjoint open strips. Let $\mathcal I$ be the set of irreducibles rooted at our fixed port and initially put $p(\gamma)=\rho^{L(\gamma)}$.

\begin{proposition}[Renewal factorization]\label{R:renew:factorization}
The masses $p(\gamma)$ sum to one. If $H,L,X,D$ denote respectively the height, visited-vertex length, signed transverse displacement and Euclidean diameter of an irreducible, then, for $s>0$ and $u\ge0$,
\[
 b(s,u):=\sum_{h\ge0}\sum_{\gamma\in\mathcal B_h}
       \rho^{L(\gamma)}e^{-sh-uL(\gamma)}
 =\frac1{1-\E_p e^{-sH-uL}}.
\]
In particular $b(s,0)\asymp s^{-3/4}$ as $s\downarrow0$. Under independent successive samples from $p$, $B_h$ is the probability that some renewal occurs at height $h$.
\end{proposition}
\begin{proof}
Unique factorization and multiplication of vertex weights give the geometric series with nonnegative terms. The finite strip estimate makes $b(s,0)$ finite for $s>0$, so $\sum p(\gamma)e^{-sH(\gamma)}<1$. Monotone convergence as $s\downarrow0$ gives $\sum p\le1$. On the other hand, summing $B_h\asymp(1+h)^{-1/4}$ gives $b(s,0)\asymp s^{-3/4}\to\infty$, forcing $\sum p=1$. Positive heights make the events that the $k$th renewal has height $h$ disjoint as $k$ varies. Summing their probabilities yields the last assertion.
\end{proof}

For a nonnegative function $f$ of an irreducible, define the height-indexed reward
\[
 p_f(j)=\E_p[f;H=j],\qquad
 B_f(h)=\sum_{\gamma\in\mathcal B_h}\rho^{L(\gamma)}
                 \sum_{I\text{ a piece of }\gamma} f(I).
\]
\begin{lemma}[Additive rewards]\label{R:renew:reward}
With convolution on nonnegative integral heights,
\[
 B_f=B*p_f*B,\qquad
 \sum_h e^{-sh}B_f(h)=b(s,0)^2\E_p[fe^{-sH}].
\]
These identities hold as identities of nonnegative sums, even when a sum is infinite. For $f=L$, $B_f(h)=M_h$.
\end{lemma}
\begin{proof}
Mark one piece in a bridge and then cut immediately before and after it. The pieces before the mark form an arbitrary bridge, the mark is one irreducible, and the pieces after it form another arbitrary bridge. This is a bijection, with multiplicative weights. Tonelli's theorem gives both formulas. Since length is additive across ports, summing the lengths of all pieces gives the full visited-vertex length.
\end{proof}

\subsection{Joint increment estimates}
The next statement concerns the joint law of the four coordinates, rather than four independent variables. Its exponent-only consequences are sufficient for the prescribed-endpoint applications.
\begin{proposition}\label{R:renew:increments}
Put $\beta=3/4$ and $\alpha=9/16$. Under $p$,
\begin{align*}
 1-\E_p e^{-sH}&\asymp s^\beta,&
 p(H>x)+p(D>x)&\le Cx^{-\beta},\\
 \E_p[L;H\le x]&\le Cx^{7/12},&
 p(L>n)&\le Cn^{-\alpha},&
 1-\E_p e^{-uL}&\asymp u^\alpha.
\end{align*}
The constants are uniform for sufficiently small positive $s,u$ and sufficiently large $x,n$.
\end{proposition}
\begin{proof}
The first assertion follows directly from Proposition~\ref{R:renew:factorization}. Since $1-e^{-H/x}\ge1-e^{-1}$ on $H>x$, it gives the height tail.

We first prove the diameter tail, which is not a consequence of the height tail alone. Set $s=A/x$ with a fixed large $A$. Give all finite bridges their normalized weight proportional to $\rho^L e^{-sH}$; its normalizer is $b(s,0)$. The mass of bridges of diameter greater than $x$ is at most
\[
 Cx^{-1/4}\sum_{h\ge1}e^{-sh}\le Cx^{-1/4}s^{-1}.
\]
Relative to $b(s,0)\ge c s^{-3/4}$ this is at most $C A^{-1/4}$. Choose $A$ so that this is at most $1/2$. Put $q=\E_p[e^{-sH};D>x]$ and $d=1-\E_p e^{-sH}$. The relative mass of bridge lists containing at least one such piece is exactly
\[
 1-\frac{d}{d+q}=\frac{q}{d+q}.
\]
A piece of diameter greater than $x$ forces the whole bridge to have that diameter. Thus $q\le d\le C_Ax^{-3/4}$. On $H\le x$, $e^{-sH}\ge e^{-A}$; on $H>x$ use the height tail. This proves $p(D>x)\le Cx^{-3/4}$.

For the reward bound, retain in Lemma~\ref{R:renew:reward} only the marked pieces of height at most $x$ and outer bridges of heights at most $x$. Their total height is at most $3x$, whence
\[
 \left(\sum_{h\le x}B_h\right)^2\E_p[L;H\le x]
 \le\sum_{h\le3x}M_h\le Cx^{25/12}.
\]
Since the squared factor on the left is at least $cx^{3/2}$, the reward is at most $Cx^{7/12}$. Taking $x=n^{3/4}$ and splitting according to $H>x$ gives
\[
 p(L>n)\le p(H>x)+n^{-1}\E_p[L;H\le x]
 \le Cn^{-9/16}.
\]
Integration of this tail yields $\E_p(1-e^{-uL})\le Cu^{9/16}$.

For the reverse bound take $s=R^{-1}$ and $u=R^{-4/3}$. On every height $R\le h\le2R$, Theorem~\ref{R:renew:finite-input} retains critical mass at least $cR^{-1/4}$ with $L\ge cR^{4/3}$. Consequently
\[
 b(s,0)-b(s,u)\ge cR^{3/4}\ge c' b(s,0).
\]
Writing $d_s=1-\E_p e^{-sH}$ and
$e_{s,u}=\E_p[e^{-sH}(1-e^{-uL})]$, the renewal identity makes the relative difference equal to $e_{s,u}/(d_s+e_{s,u})$. The displayed lower bound therefore implies $e_{s,u}\ge c d_s\ge cR^{-3/4}$. Finally $\E_p(1-e^{-uL})\ge e_{s,u}$ and $R^{-3/4}=u^{9/16}$.
\end{proof}

\begin{remark}\label{R:renew:weak-scope}
If the first-length upper bound and the mass of long bridges have only exponent precision, the same proof gives
\[
 \E_p[L;H\le x]\le x^{7/12+o(1)},\quad
 p(L>n)\le n^{-9/16+o(1)},\quad
 1-\E_p e^{-uL}=u^{9/16+o(1)}.
\]
The bounded-factor height and diameter tails are unchanged. These are the joint-law inputs used in the prescribed-endpoint arguments of Section~\ref{R:resp:section}; no endpoint conditioning is implicit in them.
\end{remark}

\subsection{Infinite concatenation and spatial growth}
We record once the elementary conversion from a Laplace deficit to the growth of independent sums.
\begin{lemma}\label{R:renew:growth}
Let $Y_1,Y_2,\ldots$ be independent copies of a nonnegative random variable $Y$, and let $0<a<1$. If $\Prob(Y>t)\le t^{-a+o(1)}$, then almost surely
$\sum_{j\le k}Y_j\le k^{1/a+o(1)}$. If also
$1-\E e^{-uY}=u^{a+o(1)}$, then almost surely
$\sum_{j\le k}Y_j=k^{1/a+o(1)}$.
\end{lemma}
\begin{proof}
Fix $\epsilon>0$ and let $k=2^j$. Put $t=k^{1/a+\epsilon}$. The probability of a summand above $t$ is at most $k t^{-a+o(1)}$. Moreover
$\E[Y\ind_{Y\le t}]\le t^{1-a+o(1)}$, by integrating the tail, so Markov's inequality bounds the probability that the truncated sum exceeds $t$ by the same expression. This is a summable geometric sequence in $j$. For the lower bound put $t=k^{1/a-\epsilon}$ and $u=t^{-1}$. Exponential Markov gives
\[
 \Prob\{\textstyle\sum_{i\le k}Y_i\le t\}
 \le e\big(\E e^{-Y/t}\big)^k
 \le\exp\{1-k t^{-a+o(1)}\},
\]
also summable. Borel--Cantelli, monotonicity between consecutive dyadic $k$, and a countable sequence of $\epsilon\downarrow0$ prove the claims.
\end{proof}

\begin{theorem}[Infinite spatial law]\label{R:renew:infinite}
Concatenate independent irreducibles of law $p$, and enumerate the resulting vertices as $\gamma_1,\gamma_2,\ldots$. Almost surely,
\[
 |\gamma_n-\gamma_1|=n^{3/4+o(1)},\qquad
 \#\{n:|\gamma_n-\gamma_1|\le R\}=R^{4/3+o(1)}.
\]
\end{theorem}
\begin{proof}
Lemma~\ref{R:renew:growth} and Proposition~\ref{R:renew:increments} give
\[
 \sum_{i\le k}L_i=k^{16/9+o(1)},\qquad
 \sum_{i\le k}H_i=k^{4/3+o(1)},\qquad
 \sum_{i\le k}D_i\le k^{4/3+o(1)}
\]
on one event of probability one. If vertex $n$ lies in piece $k+1$, it is above the last completed renewal boundary and within the total diameter of the first $k+1$ pieces. The first of the three bounds implies $k=n^{9/16+o(1)}$: both consecutive partial sums have the same exponent, so an exceptionally long current piece creates no missing interpolation interval. The height lower bound and diameter upper bound give the first conclusion. For each fixed $\epsilon>0$, all sufficiently large $n$ then obey $n^{3/4-\epsilon}\le |\gamma_n-\gamma_1|\le n^{3/4+\epsilon}$. Inverting these two inequalities gives the asserted count of vertices in a ball.
\end{proof}

\subsection{Uniform half-plane paths have the same prefix limit}
The fugacity cancellation will follow directly from factorization in Section~\ref{R:sec:thermal}. Uniform paths of one exact length require a separate ratio argument. The passage from a two-step counting ratio to the prefix limit adapts the argument of Lawler, Schramm and Werner~\cite[Appendix]{LSW2004} and Dyhr et al.~\cite[Section~2.3, Proposition~2.9]{Dyhr2011} to honeycomb ports; we prove the required ratio by local hexagon flips. Let $a_n$ count paths of $n$ visited vertices starting inward from our bottom port and otherwise unrestricted in the strict upper half-plane; put $a_0=1$.
\begin{theorem}\label{R:renew:uniform}
One has $a_{n+2}/a_n\to\rho^{-2}$. The uniform half-plane laws converge on each fixed initial path to the infinite law of Theorem~\ref{R:renew:infinite}.
\end{theorem}
\begin{proof}
First $a_n^{1/n}\to\rho^{-1}$. The upper bound is the honeycomb connective-constant theorem \cite{DCS2012}. For the lower bound, let $b_m$ count strict bridges of length $2m$ over all heights. Concatenation at a prescribed length is injective, so $b_{m+j}\ge b_mb_j$; also $b_1>0$. The divergence $\sum_m b_m\rho^{2m}=\sum_hB_h=\infty$ and supermultiplicativity imply the required root bound along even lengths. Length-two bridges fill the bounded remainders after repeating a fixed long block. Extending the terminal port by one vertex gives the odd-length bound.

A path is specified, up to bounded initial choices, by its left/right turn word. Runs have length at most four, since five equal turns close an elementary hexagon. Except for an exponentially negligible fraction at length $n$, there are at least $\delta n$ maximal runs of length three, for some fixed $\delta>0$. Indeed, compositions into parts $1,2,4$ have growth less than $\rho^{-1}$ because $\rho+\rho^2+\rho^4<1$. Assign a sufficiently small positive weight $v<1$ to every part of length three. The geometric composition series still has growth at most $\lambda<\rho^{-1}$; the words with at most $\delta n$ such parts number at most $C(v^{-\delta}\lambda)^n$. Choose $\delta$ small enough and use the root bound above.

Such a run follows a four-edge arc of a hexagonal face, with exterior entrance and exit edges. The one omitted face vertex cannot be used internally elsewhere, because two of its neighbors are already saturated. Discard faces incident to either endpoint and the bounded number of runs at the ends of the word. The remaining faces are wholly in the strict half-plane: a honeycomb face intersecting its boundary has at most three vertices strictly above it. Each retained four-edge arc may therefore be replaced by the complementary two-edge arc.

Call a wholly allowed face a slot if it avoids the endpoints, is traversed in one consecutive two- or four-edge arc using exterior entrance and exit edges, and has all its remaining vertices unused. A flip changes the length by two. Color faces with finitely many colors so that faces of one color have disjoint vertex sets. For a fixed color the slot set is unchanged by its flips. Each resulting equivalence class with $k$ slots has $\binom{k}{j}$ paths of lengths $n_0+2j$, where $j$ is the number of four-edge arcs.

Put $\theta=\rho^2/(1+\rho^2)$. For each color and each fixed $\epsilon>0$, paths of length $n$ for which $k\ge\epsilon n$ but $|j/k-\theta|\ge\epsilon$ have exponentially negligible proportion. To see this, weight every path by $\rho^{\text{length}}$. In each class shift $j$ towards $k\theta$ by $\lfloor\eta n\rfloor$ steps, with a sufficiently small fixed $\eta>0$, using the same sign for each of the two deviation classes. The binomial successive ratios show that the new weighted class count exceeds the old by $e^{c n}$. The target lengths are $n\pm2\lfloor\eta n\rfloor$, whose total critical counts are $e^{o(n)}$ by the root limit; the original denominator is also $e^{o(n)}$. The claim follows. Summing over colors, the fraction of four-edge arcs among all slots therefore tends in probability to $\theta$. Colors with fewer than $\epsilon n$ slots contribute at most a fixed multiple of $\epsilon n$, while the total slot count is at least $\delta n-O(1)$ with probability tending to one.

Double-count the edges representing one flip between lengths $n$ and $n+2$, weighting an edge by the reciprocal of the larger endpoint slot count. One flip changes that count by at most a fixed constant, since only neighboring faces can change their status. The mean incident weight tends to $1-\theta$ at length $n$ and to $\theta$ at length $n+2$. The incident sums are at most one, so the negligible exceptional paths do not affect the limits. Hence
$a_n(1-\theta+o(1))=a_{n+2}(\theta+o(1))$, proving the ratio limit.

Fix an ordered list of $k$ irreducibles, of total length $m$. Paths which start with this list and then remain strictly above its top boundary have exactly $a_{n-m}$ possible continuations. Here $m$ is even. The probability of this event tends to $\rho^m$, the product of the specified irreducible masses. For fixed $k$, the events over all such lists are disjoint and their limiting probabilities sum to one. Truncation to a finite set of lists therefore proves convergence on every fixed initial path.
\end{proof}
